% sections/conclusion.tex

The question of whether hematopoietic lineage commitment is instructed by cytokines or chosen stochastically has framed two incompatible experimental traditions for over three decades. This work shows that the question was built on a false dichotomy. A stochastic Petri net model of the GATA1/PU.1 circuit, operating across a range of molecular copy numbers, produces deterministic-appearing instruction at population scale and stochastic fate heterogeneity at single-cell scale from the same underlying circuit dynamics. The transition between regimes is governed by molecular copy number---a parameter varied implicitly between experimental assays but never before made explicit as the key mediator of the apparent contradiction.

At single-cell copy numbers, the model reveals a three-layer cascade in which the reception layer provides a structural but EPO-dose-insensitive 4:1 erythroid bias, the decision layer amplifies stochastic fluctuations to determine individual fate, and the execution layer establishes cytokine-dependent GATA1 dominance within 48 seconds, with the GATA1/PU.1 priming crossover at $t \approx 146$~s. A Phase~X cytokine-withdrawal experiment demonstrates that this erythroid state is not epigenetically sealed: EPO withdrawal at $t=500$~s produces complete fate reversal to the myeloid attractor by $t\approx 2000$~s, confirming that the model's erythroid attractor is maintained by ongoing EPOR signalling rather than by an intrinsic molecular lock. True irreversibility requires a chromatin remodelling layer absent from the current model, whose addition is the highest-priority structural extension. The EPO* threshold---the long-sought commitment switch---is not a cellular mechanism but a statistical property of the assay scale: it sharpens with the number of averaged cells and vanishes at single-cell resolution.

These insights translate directly into five falsifiable experimental predictions with explicit protocols, including an EPO pulse-duration experiment to map the minimum cytokine exposure required for sustained priming, a clonal N-ladder assay to demonstrate the EPO* artefact, and a CyTOF prediction that pGATA1 phosphorylation fraction outperforms total GATA1 as an early commitment marker. A pH sweep adds a sixth prediction: near the bifurcation, alkaline nuclear pH modestly increases the committed fraction ($\approx$\,2$\times$ over the physiological pH range) while acidic conditions produce marginally deeper PU.1 silencing in already-committed cells---two separable effects of the same $K_{\mathrm{inh}}(\mathrm{pH})$ kinetic parameter (Phase G-v3 and H). The drug discovery implications are equally concrete: EPO dose escalation in hyporesponsive patients is limited by decision-layer stochasticity rather than receptor availability, the execution-layer kinase and E3 ligase are untapped drug targets, and receptor stoichiometry normalization may be more effective than cytokine dose titration for lineage redirection. Nuclear pH modulates commitment probability by a factor of $\approx$\,2 over the physiological range and acts bidirectionally, offering a mechanistic target complementary to cytokine dosing but of secondary magnitude relative to the GCSF/EPO ratio.

Taken together, these results demonstrate that in silico experimentation with mechanistically detailed stochastic models generates hypotheses that cannot be derived from prior theoretical frameworks, directs experimental resources toward high-information experiments, and provides a quantitative vocabulary for debating biological questions that have resisted resolution at the purely experimental level. The present model provides the computational foundation for a new generation of single-cell experiments on hematopoietic commitment, and the Signal Hierarchical Petri Nets framework that implements it is available as an open tool for the systems biology community to apply to analogous fate decision systems.
